\chapter{命令系统 (Command System)}

\section{概述}

命令系统管理用户通过 \code{/} 前缀调用的斜杠命令。这些命令提供了直接的用户交互功能，不通过 LLM 处理。命令注册表在 \code{src/commands.ts}（约 25,000 行）中管理。

\section{命令注册}

\begin{lstlisting}[language=TypeScript,breaklines=true]
// commands.ts — 静态导入
import commit from './commands/commit.js'
import compact from './commands/compact/index.js'
import config from './commands/config/index.js'
import doctor from './commands/doctor/index.js'
import memory from './commands/memory/index.js'
import review, { ultrareview } from './commands/review.js'
import vim from './commands/vim/index.js'
// ... 50+ 命令导入

// 条件导入（Feature Flag 控制）
const voiceCommand = feature('VOICE_MODE')
  ? require('./commands/voice/index.js').default
  : null

const proactive = feature('PROACTIVE') || feature('KAIROS')
  ? require('./commands/proactive.js').default
  : null

const bridge = feature('BRIDGE_MODE')
  ? require('./commands/bridge/index.js').default
  : null

// 仅内部用户
const agentsPlatform = process.env.USER_TYPE === 'ant'
  ? require('./commands/agents-platform/index.js').default
  : null
\end{lstlisting}

\section{命令分类}

\subsection{Git 操作}

\begin{longtable}{ll}
\toprule
命令 & 功能 \\
\midrule
\code{/commit} & 创建 Git 提交（自动生成提交信息） \\
\code{/commit-push-pr} & 提交、推送并创建 PR \\
\code{/diff} & 查看当前变更 \\
\code{/pr\textbackslash\{\}\_comments} & 查看 PR 评论 \\
\code{/review} & 代码审查 \\
\code{/security-review} & 安全审查 \\
\bottomrule
\end{longtable}

\subsection{会话管理}

\begin{longtable}{ll}
\toprule
命令 & 功能 \\
\midrule
\code{/compact} & 压缩对话上下文 \\
\code{/resume} & 恢复之前的会话 \\
\code{/share} & 分享会话 \\
\code{/session} & 会话管理 \\
\code{/rename} & 重命名会话 \\
\code{/clear} & 清除对话历史 \\
\bottomrule
\end{longtable}

\subsection{配置与状态}

\begin{longtable}{ll}
\toprule
命令 & 功能 \\
\midrule
\code{/config} & 设置管理 \\
\code{/memory} & 持久记忆管理 \\
\code{/doctor} & 环境诊断 \\
\code{/cost} & 查看使用费用 \\
\code{/status} & 查看状态 \\
\code{/usage} & 查看使用量 \\
\code{/theme} & 切换主题 \\
\bottomrule
\end{longtable}

\subsection{工具与扩展}

\begin{longtable}{ll}
\toprule
命令 & 功能 \\
\midrule
\code{/mcp} & MCP 服务器管理 \\
\code{/skills} & 技能管理 \\
\code{/tasks} & 任务管理 \\
\code{/init} & 初始化项目配置 \\
\code{/init-verifiers} & 初始化验证器 \\
\bottomrule
\end{longtable}

\subsection{交互模式}

\begin{longtable}{ll}
\toprule
命令 & 功能 \\
\midrule
\code{/vim} & 切换 Vim 模式 \\
\code{/keybindings} & 键绑定配置 \\
\code{/color} & 颜色设置 \\
\bottomrule
\end{longtable}

\subsection{认证}

\begin{longtable}{ll}
\toprule
命令 & 功能 \\
\midrule
\code{/login} & 登录 \\
\code{/logout} & 登出 \\
\code{/install-github-app} & 安装 GitHub App \\
\code{/install-slack-app} & 安装 Slack App \\
\bottomrule
\end{longtable}

\subsection{跨平台}

\begin{longtable}{ll}
\toprule
命令 & 功能 \\
\midrule
\code{/desktop} & 传递到桌面应用 \\
\code{/mobile} & 传递到移动应用 \\
\code{/ide} & IDE 集成 \\
\code{/teleport} & 远程传送会话 \\
\bottomrule
\end{longtable}

\subsection{Feature Flag 控制的命令}

\begin{longtable}{lll}
\toprule
命令 & Feature Flag & 功能 \\
\midrule
\code{/voice} & \code{VOICE\textbackslash\{\}\_MODE} & 语音输入 \\
\code{/proactive} & \code{PROACTIVE} / \code{KAIROS} & 主动模式 \\
\code{/bridge} & \code{BRIDGE\textbackslash\{\}\_MODE} & IDE 桥接 \\
\code{/assistant} & \code{KAIROS} & 助手模式 \\
\code{/workflows} & \code{WORKFLOW\textbackslash\{\}\_SCRIPTS} & 工作流管理 \\
\code{/force-snip} & \code{HISTORY\textbackslash\{\}\_SNIP} & 强制历史片段裁剪 \\
\bottomrule
\end{longtable}

\section{命令与工具的交叉}

命令系统与工具系统有一个有趣的交叉点——\code{getSlashCommandToolSkills()}：

\begin{lstlisting}[language=TypeScript,breaklines=true]
import { getSlashCommandToolSkills } from './commands.js'
\end{lstlisting}

某些斜杠命令可以被注册为"工具技能"，使得 LLM 不仅可以在用户主动调用时使用它们，还可以在认为合适时自主调用。

\section{远程模式命令过滤}

\begin{lstlisting}[language=TypeScript,breaklines=true]
import { filterCommandsForRemoteMode, getCommands } from './commands.js'
\end{lstlisting}

在远程模式下，某些不适用的命令（如 \code{/desktop}、\code{/vim}）会被自动过滤。

\section{命令生命周期}

\begin{lstlisting}[language=TypeScript,breaklines=true]
import { notifyCommandLifecycle } from './utils/commandLifecycle.js'

// 在查询循环完成后通知命令完成
for (const uuid of consumedCommandUuids) {
  notifyCommandLifecycle(uuid, 'completed')
}
\end{lstlisting}

每个命令有明确的生命周期追踪：启动 $\rightarrow$ 执行 $\rightarrow$ 完成/失败。这支持异步命令的状态追踪和错误报告。

\section{设计启示}

\subsection{命令 vs 工具的边界}

Claude Code 将\textbf{用户发起的操作}（命令）与\textbf{LLM 发起的操作}（工具）分离，但允许交叉。这种设计提供了清晰的心智模型：
\begin{itemize}[nosep]
  \item 命令：用户主动控制，即时反馈
  \item 工具：LLM 自主调用，异步执行
\end{itemize}

\subsection{按需加载}

只在需要时加载命令实现（通过 Feature Flag 控制），保持 CLI 的启动速度。

\bigskip\noindent\rule{\textwidth}{0.4pt}\bigskip

\textbf{下一章}：\href{./chapter-08-permissions.md}{权限与安全模型}
